% =====================================================================
% CAPÍTULO 08 — MÉTRICAS DE CALIDAD DEL ERS (SGCV-IA)
% Auditoría mediante las métricas M1–M6
% Resultados definitivos de la auditoría sobre la versión final del ERS
% =====================================================================

\chapter{Métricas de calidad}
\label{sec:metricas}

\section{Diseño del instrumento de auditoría}

La auditoría emplea seis métricas derivadas del modelo de calidad de producto de
ISO/IEC 25010 y de las características de un buen conjunto de requisitos de
ISO/IEC/IEEE 29148, aplicadas sobre la versión final del ERS del sistema SGCV-IA
(\texttt{main\_corregido.tex}). Todas las métricas se expresan como una fracción
sobre un total contable. Los conteos se obtuvieron mediante revisión directa de
las secciones del ERS: la Sección 3.2, correspondiente a las tarjetas de
requisitos funcionales; la Sección 3.6, correspondiente a la tabla consolidada
de priorización MoSCoW; la Sección 4, correspondiente al diagrama y
especificación de casos de uso; y la Tabla 29, correspondiente a la Matriz de
Trazabilidad.

Para la auditoría de M2 se contrastó la prioridad registrada en las tarjetas
individuales de los requisitos funcionales con la tabla consolidada de
priorización MoSCoW, además de revisar pares de requisitos relacionados por
contenido. La revisión identificó inconsistencias de prioridad en RF-18,
RF-19, RF-20 y RF-22. Estas inconsistencias fueron corregidas y verificadas
mediante el archivo \texttt{Trello\_Import\_SGCV\_IA.csv}.

Para M6 se consideraron los defectos identificados durante la revisión de los
25 requisitos funcionales y se verificó su corrección en la versión final del
ERS. La evidencia complementaria utilizada para comprobar la prioridad y el
caso de uso asociado a cada RF corresponde al archivo
\texttt{Trello\_Import\_SGCV\_IA.csv}.

\begin{table}[h]
\centering
\footnotesize
\caption{Conteos base de la auditoría}
\label{tab:conteos}
\begin{tabular}{|p{7.2cm}|c|p{6.0cm}|}
\hline
\textbf{Concepto} & \textbf{Valor} & \textbf{Procedencia} \\ \hline
Requisitos funcionales (RF) & 25 &
Tarjetas RF-01 a RF-24 y RF-26, Sección 3.2 \\ \hline
Requisitos no funcionales (RNF) & 16 &
Columna Métrica/Criterio de aceptación de la tabla de RNF \\ \hline
Total auditado (RF + RNF) & 41 &
Suma de los anteriores \\ \hline
Casos de uso identificados & 10 &
Diagrama de casos de uso, Sección 4.1 \\ \hline
Casos de uso especificados & 10 &
Especificaciones individuales, Sección 4.2 \\ \hline
Actores humanos (actor principal) & 5 &
Administrador, Médico veterinario, Personal administrativo, Cliente, Bodeguero \\ \hline
Requisitos con criterio comprobable & 41 &
25 criterios de aceptación de RF y 16 de RNF \\ \hline
Filas RF$\to$CU en la matriz de trazabilidad & 25 &
Tabla 29 (23 filas) más matriz alterna Ley$\to$RL$\to$RF
(2 filas: RF-24, RF-26) \\ \hline
RF verificados contra el tablero Trello & 25 &
\texttt{Trello\_Import\_SGCV\_IA.csv} \\ \hline
Defectos residuales tras la corrección & 0 &
Verificación final de D-01 y D-02 \\ \hline
\end{tabular}
\end{table}

\section{Tabla de auditoría con la aritmética de cada métrica}

\begin{longtable}{|p{3.4cm}|p{3.2cm}|c|c|c|c|}
\hline
\textbf{Métrica} &
\textbf{Fórmula} &
\textbf{Aritmética} &
\textbf{Valor} &
\textbf{Ref.} &
\textbf{Cumple} \\ \hline
\endfirsthead

\hline
\textbf{Métrica} &
\textbf{Fórmula} &
\textbf{Aritmética} &
\textbf{Valor} &
\textbf{Ref.} &
\textbf{Cumple} \\ \hline
\endhead

M1a Completitud de atributos &
requisitos con los 4 atributos / requisitos &
$25/25$ &
100\,\% &
$\geq 95\,\%$ &
Sí \\ \hline

M1b Especificación de casos de uso &
CU especificados / CU identificados &
$10/10$ &
100\,\% &
100\,\% &
Sí \\ \hline

M1c Cobertura de actores &
actores con $\geq 1$ RF / actores &
$5/5$ &
100\,\% &
100\,\% &
Sí \\ \hline

M2 Consistencia &
$1-$ conflictos / pares analizados &
$1-0/31$ &
1,00 &
$\geq 0{,}98$ &
Sí \\ \hline

M3 Verificabilidad &
requisitos con criterio comprobable / requisitos &
$41/41$ &
100\,\% &
$\geq 90\,\%$ &
Sí \\ \hline

M4ade Trazabilidad adelante (RF$\to$CU) &
RF trazados a un CU específico / RF &
$25/25$ &
100\,\% &
$\geq 90\,\%$ &
Sí \\ \hline

M4atr Trazabilidad atrás &
RF con fuente / RF &
$25/25$ &
100\,\% &
100\,\% &
Sí \\ \hline

M5 Modificabilidad &
$\sum$ requisitos afectados / muestra &
$17/5$ &
3,4 &
$\leq 3{,}0$ &
No \\ \hline

M6 Corrección &
defectos residuales / requisitos &
$0/25$ &
0,00 &
$\leq 0{,}05$ &
Sí \\ \hline

\caption{Auditoría de calidad del ERS final}
\label{tab:auditoria}
\end{longtable}

\section{Corrección del enrutamiento de requisitos a casos de uso}

La revisión de la Matriz de Trazabilidad identificó que 18 de los 23 RF
registrados directamente en dicha matriz estaban asociados al caso de uso
genérico CU-02 (\emph{Registrar Atención Clínica}). La revisión de la evidencia
registrada en el tablero \texttt{Trello\_Import\_SGCV\_IA.csv} permitió
establecer el caso de uso específico correspondiente a cada requisito.

\begin{table}[h]
\centering
\footnotesize
\caption{Enrutamiento de requisitos a casos de uso}
\label{tab:correccion-cu}
\begin{tabular}{|c|c|c|}
\hline
\textbf{RF} & \textbf{CU en la matriz original} & \textbf{CU verificado} \\ \hline
RF-01 & CU-02 & CU-01 \\ \hline
RF-03 & CU-02 & CU-03 \\ \hline
RF-04 & CU-02 & CU-02 \\ \hline
RF-05 & CU-02 & CU-04 \\ \hline
RF-06 & CU-02 & CU-04 \\ \hline
RF-07 & CU-02 & CU-07 \\ \hline
RF-08 & CU-02 & CU-07 \\ \hline
RF-09 & CU-02 & CU-09 \\ \hline
RF-10 & CU-02 & CU-09 \\ \hline
RF-11 & CU-02 & CU-10 \\ \hline
RF-12 & CU-02 & CU-10 \\ \hline
RF-13 & CU-02 & CU-08 \\ \hline
RF-14 & CU-02 & CU-08 \\ \hline
RF-15 & CU-02 & CU-02 \\ \hline
RF-17 & CU-02 & CU-06 \\ \hline
RF-18 & CU-02 & CU-06 \\ \hline
RF-19 & CU-02 & CU-06 \\ \hline
RF-21 & CU-02 & CU-02 \\ \hline
\end{tabular}
\end{table}

De los 18 RF originalmente asociados a CU-02, 15 corresponden a casos de uso
específicos distintos y tres permanecen asociados a CU-02. El resultado
establece un vínculo específico entre cada requisito funcional y su caso de
uso correspondiente.

\section{Lectura interpretativa de cada métrica}

\subsubsection*{M1 Completitud}

Los tres indicadores de completitud alcanzan 100\,\%. Los 25 RF cumplen con
los atributos establecidos en la plantilla de requisitos, los 10 casos de uso
identificados cuentan con su especificación correspondiente en la Sección 4.2
y los cinco actores humanos considerados en la auditoría participan como actor
principal en al menos un caso de uso.

Por tanto, M1 alcanza el nivel máximo establecido para la auditoría.

\subsubsection*{M2 Consistencia}

M2 se calcula como uno menos la razón entre conflictos detectados y pares
analizados. Se revisaron 25 pares correspondientes a la prioridad de cada RF
frente a la tabla consolidada de MoSCoW y seis pares adicionales de contenido
entre requisitos relacionados.

La revisión identificó inconsistencias de prioridad en RF-18, RF-19, RF-20 y
RF-22. El archivo \texttt{Trello\_Import\_SGCV\_IA.csv} registra
\emph{Should have} para estos cuatro requisitos, coincidiendo con la tabla
consolidada de MoSCoW. Esta evidencia establece la prioridad registrada en la
versión final del ERS.

Los 31 pares revisados presentan consistencia, por lo que:

\[
M2 = 1-\frac{0}{31}=1{,}00
\]

El resultado cumple el umbral establecido de $\geq 0{,}98$.

\subsubsection*{M3 Verificabilidad}

Los 41 requisitos auditados, correspondientes a 25 RF y 16 RNF, cuentan con
criterios de aceptación comprobables. Los criterios incluyen umbrales
cuantificables relacionados con tiempos de respuesta, porcentajes, plazos o
número de pasos.

En el caso de RNF-01, aunque la descripción emplea el término ``rapidez'', la
métrica asociada establece un tiempo de respuesta $\leq 2$ segundos en el
95\,\% de las búsquedas. Por tanto, el requisito dispone de un criterio
comprobable.

En consecuencia:

\[
M3=\frac{41}{41}=100\,\%
\]

\subsubsection*{M4 Trazabilidad}

La trazabilidad adelante se evalúa mediante el vínculo
RF$\rightarrow$CU. Los 25 requisitos funcionales cuentan con un caso de uso
específico asociado, por lo que la cobertura alcanza:

\[
M4ade=\frac{25}{25}=100\,\%
\]

La revisión del enrutamiento permitió establecer asociaciones específicas entre
los requisitos funcionales y sus correspondientes casos de uso. Los requisitos
RF-04, RF-15 y RF-21 mantienen su asociación con CU-02.

La trazabilidad atrás también alcanza el 100\,\%, dado que los 25 RF cuentan
con una fuente documentada. Las fuentes corresponden a evidencias de
entrevistas y, para RF-24 y RF-26, al análisis normativo asociado a la LOPDP.

Por tanto, los indicadores considerados en la auditoría presentan una
cobertura del 100\,\%.

\subsubsection*{M5 Modificabilidad}

El grafo de dependencias se construyó a partir del campo
``Dependencias'' declarado en los 25 RF. Para la muestra de los cinco
requisitos con mayor grado de entrada, la suma de requisitos afectados es 17.

Por tanto:

\[
M5=\frac{17}{5}=3{,}4
\]

El valor obtenido supera el umbral establecido de $\leq 3{,}0$, por lo que
M5 no cumple el criterio definido para esta métrica.

El mayor grado de dependencia corresponde a RF-04, asociado al registro y
actualización del historial clínico. RF-02 presenta también un grado elevado
de dependencia. Estos resultados identifican una concentración de relaciones
entre requisitos dentro de los componentes funcionales del sistema.

\subsubsection*{M6 Corrección}

M6 mide la proporción de defectos residuales respecto del total de requisitos
funcionales auditados.

Durante la revisión final se verificaron los dos defectos identificados en la
auditoría: D-01, relacionado con la prioridad de RF-18, RF-19, RF-20 y RF-22,
y D-02, relacionado con el enrutamiento de requisitos hacia casos de uso.

El archivo \texttt{Trello\_Import\_SGCV\_IA.csv} permitió verificar la
prioridad y el caso de uso correspondiente. Las correcciones se encuentran
incorporadas en la versión final del ERS.

La revisión final registra cero defectos residuales entre los 25 requisitos
funcionales:

\[
M6=\frac{0}{25}=0{,}00
\]

El resultado cumple el umbral establecido de $\leq 0{,}05$.

\section{Valoración global}

\begin{table}[h]
\centering
\footnotesize
\caption{Valoración ponderada de la auditoría}
\label{tab:valoracion}
\begin{tabular}{|p{4.2cm}|c|c|c|p{5.0cm}|}
\hline
\textbf{Métrica} &
\textbf{Peso} &
\textbf{Nivel (0--4)} &
\textbf{Aporte} &
\textbf{Criterio del nivel} \\ \hline

M1 Completitud &
0,20 &
4 &
0,80 &
Los tres indicadores alcanzan 100\,\% \\ \hline

M2 Consistencia &
0,15 &
4 &
0,60 &
1,00 sobre 31 pares revisados \\ \hline

M3 Verificabilidad &
0,20 &
4 &
0,80 &
41 de 41 requisitos con criterio comprobable \\ \hline

M4 Trazabilidad &
0,20 &
2 &
0,40 &
Los indicadores de trazabilidad considerados alcanzan 100\,\% de cobertura \\ \hline

M5 Modificabilidad &
0,10 &
2 &
0,20 &
3,4 frente al umbral establecido de $\leq 3{,}0$ \\ \hline

M6 Corrección &
0,15 &
4 &
0,60 &
0,00 defectos residuales sobre 25 RF \\ \hline

\textbf{Total} &
\textbf{1,00} &
--- &
\textbf{3,40} &
Sobre 4,00 (85\,\%) \\ \hline

\end{tabular}
\end{table}

La valoración ponderada de la auditoría es:

\begin{equation}
(4\times0{,}20)+(4\times0{,}15)+(4\times0{,}20)
+(2\times0{,}20)+(2\times0{,}10)+(4\times0{,}15)
=3{,}40
\end{equation}

El resultado corresponde a una valoración de:

\[
\boxed{3{,}40\text{ sobre }4{,}00\;(85\,\%)}
\]

La auditoría muestra cumplimiento en M1, M2, M3 y M6. M4 presenta una
cobertura del 100\,\% en los vínculos de trazabilidad evaluados, mientras que
M5 obtiene un valor de 3,4, superior al umbral de 3,0 definido para la
métrica.

En conjunto, la versión final del ERS alcanza una valoración ponderada de
3,40 sobre 4,00, equivalente al 85\,\% del nivel máximo de calidad considerado
en la auditoría.
